% Added Extension boxes for Tutorial 3 Q2 Method 2; Tutorial 3 Q4 Method 2; Tutorial 3 Q5 Method 2; Tutorial 3 Q6 Method 2; Tutorial 3 Q7 Method 2

%====================================================================

\section{Finite-Dimensional Vector Spaces}

\label{sec:finite-dimensional-vector-spaces}

%====================================================================

This chapter corresponds to Study Weeks 3--4 (ILOs 2--3). The focus is on \emph{finite-dimensional} structure: linear independence, bases, coordinates, and dimension arguments.

%--------------------------------------------------------------------

\subsection{Span and Generating Sets}

%--------------------------------------------------------------------

Throughout, let $V$ be a vector space over a field $\mathbb{F}$ (typically $\mathbb{R}$ or $\mathbb{C}$).

\begin{definition}{Linear combination and span}{ch2-lincomb-span}

Let $S \subseteq V$.

\begin{itemize}

\item A \textbf{linear combination} of vectors in $S$ means an expression

\[
\lambda_1 v_1+\cdots+\lambda_n v_n
\]

where $n\in\mathbb{N}$, $v_1,\ldots,v_n\in S$ are \emph{distinct}, and $\lambda_1,\ldots,\lambda_n\in\mathbb{F}$.

\item The \textbf{span} of $S$ is

\[
\operatorname{Span}(S)\;=\;\left\{\lambda_1 v_1+\cdots+\lambda_n v_n \,\middle|\,
n\in\mathbb{N},\ v_1,\ldots,v_n\in S\text{ distinct},\ \lambda_i\in\mathbb{F}\right\}.
\]

\item We say $S$ \textbf{generates} (or \textbf{spans}) $V$ if $\operatorname{Span}(S)=V$.

\end{itemize}

\end{definition}

\begin{theorem}{Span is the smallest subspace containing a set}{ch2-span-smallest}

Let $S\subseteq V$.

\begin{enumerate}

\item $\operatorname{Span}(S)$ is a vector subspace of $V$ and contains $S$.

\item If $W$ is a subspace of $V$ and $S\subseteq W$, then $\operatorname{Span}(S)\subseteq W$.

\end{enumerate}

Hence $\operatorname{Span}(S)$ is the \emph{smallest} subspace of $V$ containing $S$.

\begin{proof}

(1) By definition, $0\in\operatorname{Span}(S)$ (take all coefficients $0$), and the set is closed under addition and scalar multiplication because sums/scalar multiples of finite linear combinations are again finite linear combinations of elements of $S$. Also $S\subseteq \operatorname{Span}(S)$ since each $s\in S$ equals $1\cdot s$.

(2) If $S\subseteq W$ and $W$ is a subspace, then every linear combination of vectors in $S$ lies in $W$ by closure. Therefore $\operatorname{Span}(S)\subseteq W$.

\end{proof}

\end{theorem}

\begin{commonmistake}

\begin{itemize}

\item Confusing $S$ with $\operatorname{Span}(S)$: the span is usually \emph{much larger} than the set.

\item Forgetting that $\operatorname{Span}(\varnothing)=\{0\}$ (the empty linear combination is $0$).

\end{itemize}

\end{commonmistake}

%--------------------------------------------------------------------

\subsection{Linear Independence and Dependence}

%--------------------------------------------------------------------

\begin{definition}{Linear independence}{ch2-linear-independence}

Let $S\subseteq V$ (not necessarily finite).

\begin{itemize}

\item $S$ is \textbf{linearly independent} if whenever

\[
\lambda_1 v_1+\cdots+\lambda_n v_n = 0
\]

for some $n\in\mathbb{N}$, distinct $v_1,\ldots,v_n\in S$, and scalars $\lambda_i\in\mathbb{F}$, then

\[
\lambda_1=\cdots=\lambda_n=0.
\]

\item Otherwise, $S$ is \textbf{linearly dependent}.

\end{itemize}

\end{definition}

% --- ADDED: Intuition for linear independence ---
Intuitively, a set $S$ is linearly independent when \emph{no vector in $S$ is redundant}: you cannot ``build'' any element of $S$ from the others.
Equivalently, the only way to combine vectors from $S$ to get $0$ is by using all-zero coefficients---there is no ``shortcut'' or ``hidden relation.''
The negation (linear dependence) means some vector in $S$ \emph{can} be expressed using the others; this is made precise in Theorem~\ref{thm:ch2-dep-iff-in-span}.

% --- ADDED: Quick facts about subsets ---
\begin{policynote}
\textbf{Useful quick facts about linear independence.}
\begin{itemize}[leftmargin=*]
\item \emph{Any subset of a linearly independent set is linearly independent.}  (If a finite sub-selection could give a nontrivial relation, that same relation would work in the larger set.)
\item \emph{Any superset of a linearly dependent set is linearly dependent.}  (A nontrivial relation among a subset extends to the whole set by assigning coefficient $0$ to the extra vectors.)
\item $\varnothing$ is linearly independent (vacuously: there are no finite selections to test).
\item $\{v\}$ is linearly independent if and only if $v\neq 0$.
\item $\{u,v\}$ is linearly dependent if and only if one is a scalar multiple of the other.
\end{itemize}
\end{policynote}

\begin{keyformula}

\textbf{Finite-set test (most used in computations).}

For distinct $v_1,\ldots,v_n\in V$,

\[
\begin{aligned}
\{v_1,\ldots,v_n\}\text{ is linearly independent}
\quad\Longleftrightarrow\quad&
\lambda_1 v_1+\cdots+\lambda_n v_n=0\\
&\Rightarrow\ \lambda_1=\cdots=\lambda_n=0.
\end{aligned}
\]

In $\mathbb{F}^m$, place $v_1,\ldots,v_n$ as columns of a matrix $A$:

\[
A=[v_1\ \cdots\ v_n].
\]

Then $\{v_1,\ldots,v_n\}$ is linearly independent $\Longleftrightarrow$ the homogeneous system $A\lambda=0$ has only the trivial solution $\Longleftrightarrow$ every column is a pivot column in $\operatorname{rref}(A)$.

\end{keyformula}

% --- ADDED: Dimension sanity check ---
\begin{examtip}
\textbf{Quick dimension sanity check for independence.}
In $\mathbb{F}^m$, any set of more than $m$ vectors is automatically linearly dependent (you have more unknowns than equations in the homogeneous system, so a nontrivial solution must exist).
In particular, in $\mathbb{R}^{2,2}\cong\mathbb{R}^4$, at most $4$ matrices can be independent; in $P_n(\mathbb{F})$, at most $n+1$ polynomials can be independent.
\end{examtip}

\begin{theorem}{Dependence iff one vector lies in the span of the others (Tutorial 3, Question 1)}{ch2-dep-iff-in-span}

Let $S\subseteq V$. The following are equivalent:

\begin{enumerate}

\item $S$ is linearly dependent.

\item There exist $n\in\mathbb{N}\cup\{0\}$ and $n+1$ distinct vectors $v,w_1,\ldots,w_n\in S$ such that

\[
v\in \operatorname{Span}(\{w_1,\ldots,w_n\}).
\]

(For $n=0$, this means $v\in\operatorname{Span}(\varnothing)=\{0\}$, i.e.\ $v=0$.)

\end{enumerate}

\begin{proof}

$(1)\Rightarrow(2)$: If $S$ is dependent, then there exist distinct $s_0,\ldots,s_m\in S$ and scalars $\alpha_0,\ldots,\alpha_m$, not all zero, such that

\[
\alpha_0 s_0+\cdots+\alpha_m s_m=0.
\]

Choose $k$ with $\alpha_k\neq 0$ and solve for $s_k$:

\[
s_k=-\sum_{i\neq k}\frac{\alpha_i}{\alpha_k}s_i\in \operatorname{Span}(\{s_0,\ldots,s_m\}\setminus\{s_k\}).
\]

Set $v=s_k$ and list the remaining vectors as $w_1,\ldots,w_n$.

$(2)\Rightarrow(1)$: If $v=\beta_1 w_1+\cdots+\beta_n w_n$, then

\[
v-\beta_1 w_1-\cdots-\beta_n w_n=0
\]

is a nontrivial linear dependence relation among distinct vectors from $S$ (the coefficient of $v$ is $1\neq 0$). Hence $S$ is dependent.

\end{proof}

\end{theorem}

% --- ADDED: Proof note on the minimal-dependent-subset refinement ---
\begin{policynote}
\textbf{Refinement via minimal dependent subsets (Tutorial 3, Question 1, Method 2).}
One can strengthen the $(1)\Rightarrow(2)$ direction by choosing a \emph{minimal} finite dependent subset $T\subseteq S$ (one of smallest cardinality).  In such a minimal set, \emph{every} coefficient in the dependence relation is nonzero (if some $\alpha_k=0$, the smaller set $T\setminus\{s_k\}$ would still be dependent, contradicting minimality).
This is useful in proofs where you need to remove a \emph{specific} vector (e.g.\ the one with the largest index, or one satisfying an extra condition) and be certain its coefficient is nonzero.
\end{policynote}

\begin{examplebox}

\textbf{Worked Example (Week 3 lecture): linear independence in $\mathbb{R}^{2,2}$.}

Determine whether

\[
M_1=\begin{pmatrix}1&0\\-2&1\end{pmatrix},\
M_2=\begin{pmatrix}0&-1\\1&1\end{pmatrix},\
M_3=\begin{pmatrix}-1&2\\1&0\end{pmatrix},\
M_4=\begin{pmatrix}2&1\\2&-2\end{pmatrix}
\]

is a linearly independent subset of $\mathbb{R}^{2,2}$.

\medskip

\textbf{Solution.}

Consider

\[
c_1M_1+c_2M_2+c_3M_3+c_4M_4=0.
\]

Equating entries gives a homogeneous system in $(c_1,c_2,c_3,c_4)$:

\[
\begin{pmatrix}
1 & 0 & -1 & 2\\
0 & -1 & 2 & 1\\
-2 & 1 & 1 & 2\\
1 & 1 & 0 & -2
\end{pmatrix}
\begin{pmatrix}c_1\\c_2\\c_3\\c_4\end{pmatrix}
=
\begin{pmatrix}0\\0\\0\\0\end{pmatrix}.
\]

Row reduction yields $\operatorname{rref}=I_4$, so the only solution is $c_1=c_2=c_3=c_4=0$.

Hence $\{M_1,M_2,M_3,M_4\}$ is linearly independent.

\end{examplebox}

% --- ADDED: General method note for non-R^n spaces ---
\begin{policynote}
\textbf{Strategy for non-$\mathbb{R}^n$ spaces (matrices, polynomials, functions).}
The worked example above illustrates the general recipe:
\begin{enumerate}[leftmargin=*]
\item \emph{Flatten} the vectors into coordinate vectors (read off entries of matrices, or coefficients of polynomials).
\item \emph{Form a matrix} with those coordinate vectors as columns.
\item \emph{Row reduce} to determine whether every column is a pivot column.
\end{enumerate}
For function spaces (e.g.\ $\{e^x,e^{-x}\}$ in $\mathbb{R}^{\mathbb{R}}$), there is no natural finite coordinate system, so instead you \emph{evaluate at specific points} to generate equations, as in the next example.
\end{policynote}

\begin{examplebox}

\textbf{Worked Example (Week 3 lecture): linear independence of functions.}

Let $f,g:\mathbb{R}\to\mathbb{R}$ be defined by $f(x)=e^x$ and $g(x)=e^{-x}$. Show that $\{f,g\}$ is linearly independent.

\medskip

\textbf{Solution.}

Assume $\alpha f+\beta g=0$ as functions. Then for all $x\in\mathbb{R}$,

\[
\alpha e^x+\beta e^{-x}=0.
\]

Multiply by $e^x$:

\[
\alpha e^{2x}+\beta=0 \quad\text{for all }x.
\]

Setting $x=0$ gives $\alpha+\beta=0$, and setting $x=1$ gives $\alpha e^{2}+\beta=0$.

Subtracting the two equations yields $\alpha(e^2-1)=0$, so $\alpha=0$, hence $\beta=0$.

Therefore $\{f,g\}$ is linearly independent.

\end{examplebox}

% --- ADDED: Lecture Week 3 example on {v1-v2, v2-v3, v3-v4, v4} ---
\begin{examplebox}
\textbf{Worked Example (Week 3 lecture): independence of differences.}

Let $v_1,v_2,v_3,v_4$ be distinct vectors such that $\{v_1,v_2,v_3,v_4\}$ is linearly independent.
Show that $\{v_1-v_2,\;v_2-v_3,\;v_3-v_4,\;v_4\}$ is linearly independent.

\medskip

\textbf{Solution.}
Suppose $\alpha_1(v_1-v_2)+\alpha_2(v_2-v_3)+\alpha_3(v_3-v_4)+\alpha_4 v_4=0$.
Expanding:
\[
\alpha_1 v_1+(\alpha_2-\alpha_1)v_2+(\alpha_3-\alpha_2)v_3+(\alpha_4-\alpha_3)v_4=0.
\]
By independence of $\{v_1,v_2,v_3,v_4\}$, all coefficients are zero:
\[
\alpha_1=0,\quad \alpha_2-\alpha_1=0,\quad \alpha_3-\alpha_2=0,\quad \alpha_4-\alpha_3=0.
\]
This gives $\alpha_1=\alpha_2=\alpha_3=\alpha_4=0$.
Hence $\{v_1-v_2,\;v_2-v_3,\;v_3-v_4,\;v_4\}$ is linearly independent.

\medskip
\emph{Key idea:} express the new vectors in terms of the old basis and check that the ``change-of-basis'' matrix is invertible.  Here the matrix is lower-triangular with $1$'s on the diagonal, hence invertible.
\end{examplebox}

\begin{commonmistake}

\begin{itemize}

\item Using the finite-set definition incorrectly for infinite sets: linear dependence is witnessed by a \emph{finite} nontrivial linear combination.

\item Forgetting "distinct vectors" in the definition: repeated vectors can hide dependence in notation.

\end{itemize}

\end{commonmistake}

%--------------------------------------------------------------------

\subsection{Bases and Coordinate Representations}

%--------------------------------------------------------------------

\begin{definition}{Basis}{ch2-basis}

A \textbf{basis} for $V$ is a subset $B\subseteq V$ such that

\begin{enumerate}

\item $B$ is linearly independent;

\item $\operatorname{Span}(B)=V$.

\end{enumerate}

If $B$ is finite, we often write it as an ordered list $(b_1,\ldots,b_n)$ when computing coordinates.

\end{definition}

% --- ADDED: Intuition for a basis ---
A basis is a ``coordinate system'' for $V$: it provides a minimal set of building blocks from which every vector can be uniquely assembled.
Condition~(1) (independence) ensures no building block is redundant; condition~(2) (spanning) ensures every vector is reachable.
Together they give the key consequence: every $v\in V$ has a \emph{unique} representation as a linear combination of basis vectors.

% --- ADDED: Standard bases ---
\begin{keyformula}
\textbf{Standard bases (know these by heart).}
\begin{itemize}[leftmargin=*]
\item $\mathbb{F}^n$: the standard basis is $\{e_1,\ldots,e_n\}$ where $e_i$ has a $1$ in position $i$ and $0$'s elsewhere.
\item $P_n(\mathbb{F})$: the monomial basis is $\{1,X,X^2,\ldots,X^n\}$; note $\dim(P_n(\mathbb{F}))=n+1$.
\item $\mathbb{F}^{m,n}$: the matrix-unit basis is $\{E_{ij}\mid 1\le i\le m,\,1\le j\le n\}$; note $\dim(\mathbb{F}^{m,n})=mn$.
\item $P(\mathbb{F})$ (all polynomials): the infinite monomial basis $\{1,X,X^2,\ldots\}$; $\dim(P(\mathbb{F}))=\infty$.
\end{itemize}
\end{keyformula}

\begin{theorem}{Unique linear combination in a basis}{ch2-unique-lincomb-basis}

Let $B$ be a basis of $V$. Then each $v\in V$ can be written as a linear combination of vectors of $B$ in \emph{exactly one} way (as a finite sum):

\[
v=\sum_{b\in B}\lambda_b\,b,
\]

where only finitely many $\lambda_b\in\mathbb{F}$ are nonzero.

\begin{proof}

Existence: since $\operatorname{Span}(B)=V$, every $v$ is a finite linear combination of elements of $B$.

Uniqueness: if also $v=\sum_{b\in B}\mu_b\,b$ (finite sums), subtracting gives

\[
\sum_{b\in B}(\lambda_b-\mu_b)b=0
\]

as a finite linear combination of distinct vectors in $B$. Linear independence forces $\lambda_b-\mu_b=0$ for all $b$, so $\lambda_b=\mu_b$.

\end{proof}

\end{theorem}

% --- ADDED: Basis Existence Theorem ---
\begin{theorem}{Basis Existence Theorem (Lecture Week 3)}{ch2-basis-existence}
Every vector space has a basis.
\end{theorem}

This theorem is equivalent to the Axiom of Choice and is a pure existence result.
In practice, for the spaces encountered in MH1201 ($\mathbb{F}^n$, $P_n(\mathbb{F})$, $\mathbb{F}^{m,n}$, etc.), explicit bases are readily constructed.
There exist vector spaces (e.g.\ $\mathbb{R}^{\mathbb{R}}$) for which no explicit basis can be written down, even though one must exist.

\begin{definition}{Coordinate vector (finite basis)}{ch2-coordinate-vector}

Let $V$ be a vector space with an \emph{ordered} basis $\mathcal{B}=(b_1,\ldots,b_n)$.

For $v\in V$, write uniquely

\[
v=\lambda_1 b_1+\cdots+\lambda_n b_n.
\]

The column vector

\[
[v]_{\mathcal{B}}\;\stackrel{\mathrm{def}}{=}\;
\begin{pmatrix}\lambda_1\\ \vdots\\ \lambda_n\end{pmatrix}\in\mathbb{F}^{n,1}
\]

is called the \textbf{coordinate vector} of $v$ with respect to $\mathcal{B}$.

\end{definition}

% --- ADDED: Ordering matters ---
\begin{commonmistake}
\textbf{Ordering matters for coordinate vectors.}
The coordinate vector depends on the \emph{order} of the basis vectors.
If $\mathcal{B}=(b_1,b_2)$ and $\mathcal{B}'=(b_2,b_1)$, and $v=3b_1+5b_2$, then
\[
[v]_{\mathcal{B}}=\begin{pmatrix}3\\5\end{pmatrix},
\qquad
[v]_{\mathcal{B}'}=\begin{pmatrix}5\\3\end{pmatrix}.
\]
Always check which ordering the problem specifies.
\end{commonmistake}

\begin{examplebox}

\textbf{Worked Example: compute coordinates in a nonstandard basis.}

Let $\mathcal{B}=(b_1,b_2,b_3)$ be a basis of $\mathbb{R}^3$ where

\[
b_1=\begin{pmatrix}1\\0\\1\end{pmatrix},\quad
b_2=\begin{pmatrix}0\\1\\1\end{pmatrix},\quad
b_3=\begin{pmatrix}1\\1\\0\end{pmatrix}.
\]

Find $[v]_{\mathcal{B}}$ for $v=\begin{pmatrix}2\\1\\3\end{pmatrix}$.

\medskip

\textbf{Solution.}

Solve $v=\lambda_1 b_1+\lambda_2 b_2+\lambda_3 b_3$:

\[
\begin{pmatrix}2\\1\\3\end{pmatrix}
=
\lambda_1\begin{pmatrix}1\\0\\1\end{pmatrix}
+\lambda_2\begin{pmatrix}0\\1\\1\end{pmatrix}
+\lambda_3\begin{pmatrix}1\\1\\0\end{pmatrix}
=
\begin{pmatrix}\lambda_1+\lambda_3\\ \lambda_2+\lambda_3\\ \lambda_1+\lambda_2\end{pmatrix}.
\]

Thus

\[
\lambda_1+\lambda_3=2,\quad \lambda_2+\lambda_3=1,\quad \lambda_1+\lambda_2=3.
\]

Subtract the first two equations from the third:

\[
(\lambda_1+\lambda_2)-(\lambda_1+\lambda_3)-(\lambda_2+\lambda_3)=3-2-1 \Rightarrow -2\lambda_3=0 \Rightarrow \lambda_3=0.
\]

Then $\lambda_1=2$ and $\lambda_2=1$. Hence

\[
[v]_{\mathcal{B}}=\begin{pmatrix}2\\1\\0\end{pmatrix}.
\]

\end{examplebox}

% --- ADDED: General method for coordinate computations ---
\begin{policynote}
\textbf{General method for finding $[v]_{\mathcal{B}}$ in $\mathbb{R}^n$.}
Form the augmented matrix $[b_1\ b_2\ \cdots\ b_n\ |\ v]$ and row reduce.
The solution gives the coordinates $\lambda_1,\ldots,\lambda_n$.
Equivalently, if $P$ is the matrix with columns $b_1,\ldots,b_n$, then $[v]_{\mathcal{B}}=P^{-1}v$.
\end{policynote}

\begin{examtip}

When asked to show a set $B$ is a basis, do \emph{not} argue informally. Use one of:

\begin{itemize}

\item show \textbf{spanning} + \textbf{independence}; or

\item in finite dimensions, show $|B|=\dim(V)$ and either spanning or independence (then the other follows by theorem).

\end{itemize}

\end{examtip}

%--------------------------------------------------------------------

\subsection{Dimension and Core Theorems in Finite Dimensions}

%--------------------------------------------------------------------

\begin{theorem}{Size-Limitation Lemma}{ch2-size-limitation}

Let $V$ be a vector space that has a \emph{finite} basis of size $n$.

Then every linearly independent subset of $V$ is finite and has size $\le n$.

\begin{proof}

Let $B=\{b_1,\ldots,b_n\}$ be a basis of $V$. Suppose for contradiction that $V$ has a linearly independent set containing $n+1$ distinct vectors $v_1,\ldots,v_{n+1}$.

Each $v_i$ is a linear combination of $b_1,\ldots,b_n$, so there exist scalars $\lambda_{i,1},\ldots,\lambda_{i,n}\in\mathbb{F}$ with

\[
v_i=\lambda_{i,1}b_1+\cdots+\lambda_{i,n}b_n \qquad (i=1,\ldots,n+1).
\]

For scalars $\mu_1,\ldots,\mu_{n+1}$, we have

\[
\mu_1v_1+\cdots+\mu_{n+1}v_{n+1}
=
\sum_{j=1}^{n}\left(\lambda_{1,j}\mu_1+\cdots+\lambda_{n+1,j}\mu_{n+1}\right)b_j.
\]

Thus $\mu_1v_1+\cdots+\mu_{n+1}v_{n+1}=0$ holds whenever

\[
\lambda_{1,j}\mu_1+\cdots+\lambda_{n+1,j}\mu_{n+1}=0 \qquad (j=1,\ldots,n).
\]

This is a homogeneous system of $n$ equations in $n+1$ unknowns, so it has a nontrivial solution $(\mu_1,\ldots,\mu_{n+1})\neq 0$.

Hence there is a nontrivial linear combination of $v_1,\ldots,v_{n+1}$ equal to $0$, contradicting linear independence.

\end{proof}

\end{theorem}

% --- ADDED: Proof strategy note ---
\begin{policynote}
\textbf{Why does the Size-Limitation Lemma work?}
The key idea is a \emph{counting argument}: $n+1$ vectors, each expressed in terms of $n$ basis vectors, give an under-determined homogeneous system ($n$ equations, $n+1$ unknowns).
Under-determined homogeneous systems always have nontrivial solutions (this is a fundamental fact from MH1200 / linear systems theory).
This lemma is the engine behind almost every dimension result in the course.
\end{policynote}

\begin{theorem}{Fundamental Theorem of MH1201 (basis size is well-defined for finite bases)}{ch2-fundamental-mh1201}

If $V$ has a finite basis of size $n$, then every basis of $V$ is finite and has size $n$.

\begin{proof}

Let $B$ be a basis with $|B|=n$, and let $\widetilde{B}$ be any other basis.

Since $\widetilde{B}$ is linearly independent, the Size-Limitation Lemma gives $|\widetilde{B}|\le n$.

Applying the same lemma with $\widetilde{B}$ as the finite basis yields $n\le |\widetilde{B}|$.

Therefore $|\widetilde{B}|=n$.

\end{proof}

\end{theorem}

% --- ADDED: Significance note ---
\begin{policynote}
\textbf{Why ``Fundamental''?}
This theorem guarantees that dimension is an intrinsic property of the vector space, not an artefact of a particular basis choice.
Without it, we could not define $\dim(V)$ unambiguously.
The proof uses the Size-Limitation Lemma \emph{twice}, with the two bases swapping roles---this ``ping-pong'' argument is a standard proof technique.
\end{policynote}

\begin{definition}{Dimension and finite-dimensionality}{ch2-dimension}

Let $V$ be a vector space. If $V$ has a finite basis, define

\[
\dim(V)\;\stackrel{\mathrm{def}}{=}\;\text{the number of vectors in any (hence every) basis of }V.
\]

Then $V$ is \textbf{finite-dimensional} iff $\dim(V)$ is finite.

\end{definition}

% --- ADDED: Quick reference table ---
\begin{keyformula}
\textbf{Quick reference: dimensions of common spaces.}

\medskip

\begin{tabular}{ll}
Space & Dimension \\
\hline
$\mathbb{F}^n$ & $n$ \\
$P_n(\mathbb{F})$ & $n+1$ \\
$\mathbb{F}^{m,n}$ & $mn$ \\
$\{M\in\mathbb{F}^{n,n}\mid \operatorname{trace}(M)=0\}$ & $n^2-1$ \\
$\{M\in\mathbb{R}^{n,n}\mid M^T=M\}$ (symmetric) & $n(n+1)/2$ \\
$\{M\in\mathbb{R}^{n,n}\mid M^T=-M\}$ (skew-symmetric) & $n(n-1)/2$ \\
$\{M\in\mathbb{R}^{n,n}\mid M\text{ upper triangular}\}$ & $n(n+1)/2$ \\
\end{tabular}
\end{keyformula}

\begin{keyformula}

\textbf{Equivalences in finite dimensions.}

Let $\dim(V)=n$ and let $S\subseteq V$ be a set with exactly $n$ vectors. Then the following are equivalent:

\[
S\text{ is a basis}\quad\Longleftrightarrow\quad S\text{ spans }V \quad\Longleftrightarrow\quad S\text{ is linearly independent}.
\]

Also:

\begin{itemize}

\item any linearly independent set has size $\le n$;

\item any spanning set has size $\ge n$.

\end{itemize}

\end{keyformula}

% --- ADDED: Why this equivalence is so powerful ---
\begin{policynote}
\textbf{Why is this equivalence so useful in practice?}
It halves your work.  If you know $|S|=\dim(V)$, you only need to prove \emph{one} of the two conditions (independence or spanning); the other is then automatic.
This is the go-to shortcut on exams: count vectors, compare with the known dimension, and check whichever condition is easier.
\end{policynote}

\begin{extension}[title={Extension (Tutorial 3, Question 2, Method 2)}]

Let $S=\{s_1,\ldots,s_n\}\subseteq V$ and define $T:\mathbb{F}^n\to V$ by

$T(\alpha_1,\ldots,\alpha_n)=\sum_{i=1}^n \alpha_i s_i$.

Then $\operatorname{Im}(T)=\operatorname{Span}(S)$ and $\ker(T)$ is the space of linear relations among the $s_i$.

If $n=\dim(V)$ and $S$ is linearly independent, then $\ker(T)=\{0\}$ so $T$ is injective; by rank--nullity, $\dim(\operatorname{Im}(T))=n=\dim(V)$, hence $\operatorname{Im}(T)=V$ and $S$ spans $V$.

\end{extension}

\begin{examplebox}

\textbf{Worked Example (Week 3 lecture): basis for a trace-zero subspace.}

Let

\[
W=\left\{M\in\mathbb{R}^{2,2}\ \middle|\ \operatorname{trace}(M)=0\right\}.
\]

Find a basis of $W$ and compute $\dim(W)$.

\medskip

\textbf{Solution.}

Write $M=\begin{pmatrix}a&b\\c&d\end{pmatrix}$. The constraint $\operatorname{trace}(M)=a+d=0$ implies $d=-a$, so

\[
M=\begin{pmatrix}a&b\\c&-a\end{pmatrix}
=
a\begin{pmatrix}1&0\\0&-1\end{pmatrix}
+b\begin{pmatrix}0&1\\0&0\end{pmatrix}
+c\begin{pmatrix}0&0\\1&0\end{pmatrix}.
\]

The three displayed matrices are linearly independent (distinct support patterns in entries), hence they form a basis.

Therefore $\dim(W)=3$.

\end{examplebox}

% --- ADDED: Kernel viewpoint for trace-zero ---
\begin{policynote}
\textbf{Kernel viewpoint for the trace-zero subspace.}
The trace map $\operatorname{tr}:\mathbb{R}^{2,2}\to\mathbb{R}$ is a linear functional (a linear map to the scalar field).
Since $\operatorname{tr}\neq 0$ (e.g.\ $\operatorname{tr}(I)=2$), it is surjective, so $\operatorname{rank}(\operatorname{tr})=1$.
By rank--nullity, $\dim(\ker(\operatorname{tr}))=\dim(\mathbb{R}^{2,2})-1=4-1=3$.
This gives $\dim(W)=3$ without explicitly finding a basis---useful as a sanity check.
\end{policynote}

%--------------------------------------------------------------------

\subsection{Adding Vectors, Extending Bases, and Extracting Bases}

%--------------------------------------------------------------------

\begin{theorem}{Add-One-More Lemma}{ch2-add-one-more}

Let $S\subseteq V$ be linearly independent and let $v\in V$.

Then the following are equivalent:

\begin{enumerate}

\item $v\notin S$ and $S\cup\{v\}$ is linearly independent.

\item $v\notin \operatorname{Span}(S)$.

\end{enumerate}

\begin{proof}

$(1)\Rightarrow(2)$: If $v\in \operatorname{Span}(S)$, then $v=\sum_{i=1}^n\lambda_i s_i$ for distinct $s_i\in S$, so

\[
1\cdot v-\sum_{i=1}^n\lambda_i s_i=0
\]

is a nontrivial dependence among distinct vectors in $S\cup\{v\}$, contradiction.

$(2)\Rightarrow(1)$: If $v\notin \operatorname{Span}(S)$, then $v\notin S$ since $S\subseteq \operatorname{Span}(S)$. If a nontrivial linear combination of distinct vectors in $S\cup\{v\}$ equals $0$, the coefficient of $v$ must be nonzero (otherwise we contradict independence of $S$), and then we can solve for $v$ as a linear combination of vectors in $S$, contradicting $v\notin\operatorname{Span}(S)$.

\end{proof}

\end{theorem}

% --- ADDED: Why this lemma matters ---
\begin{policynote}
\textbf{The Add-One-More Lemma in action.}
This lemma is the workhorse behind \emph{both} the Basis Extension Theorem (keep adding vectors outside the span until you reach a basis) and the hereditary dimension theorem (a maximal independent set in a subspace must span it, because otherwise we could add one more).
It encodes the precise boundary between ``can still grow'' and ``already spans.''
\end{policynote}

\begin{theorem}{Finite-dimensionality is hereditary}{ch2-hereditary}

Let $V$ be finite-dimensional and let $W$ be a subspace of $V$. Then $W$ is finite-dimensional.

\begin{proof}

Let $n=\dim(V)$. Any linearly independent subset of $W$ is also linearly independent in $V$, hence has size $\le n$ by the Size-Limitation Lemma. Take a linearly independent subset $B\subseteq W$ of maximal size (existence follows since sizes are bounded by $n$). We claim $\operatorname{Span}(B)=W$. If not, pick $v\in W\setminus \operatorname{Span}(B)$; by the Add-One-More Lemma, $B\cup\{v\}$ is linearly independent in $W$, contradicting maximality. Hence $B$ is a finite basis of $W$.

\end{proof}

\end{theorem}

% --- ADDED: Corollary on dimension inequality ---
\begin{keyformula}
\textbf{Dimension inequality for subspaces.}
If $W$ is a subspace of a finite-dimensional space $V$, then:
\[
\dim(W)\le\dim(V),
\]
with equality if and only if $W=V$.

\emph{Proof of equality case:} if $\dim(W)=\dim(V)=n$, any basis of $W$ is an independent set of $n$ vectors in $V$, hence a basis of $V$ by the equivalence theorem, so $W=V$.
\end{keyformula}

\begin{theorem}{Basis Extension Theorem}{ch2-basis-extension}

Let $V$ be finite-dimensional and let $S\subseteq V$ be linearly independent. Then $S$ can be extended to a basis of $V$.

\begin{proof}

Let $n=\dim(V)$. Among all linearly independent subsets of $V$ that contain $S$, pick one of maximal size, call it $B$ (sizes are bounded by $n$). If $\operatorname{Span}(B)\neq V$, choose $v\in V\setminus\operatorname{Span}(B)$; then $B\cup\{v\}$ is linearly independent by the Add-One-More Lemma, contradicting maximality. Hence $\operatorname{Span}(B)=V$, so $B$ is a basis extending $S$.

\end{proof}

\end{theorem}

% --- ADDED: Practical algorithm from Week 4 lecture ---
\begin{policynote}
\textbf{Practical algorithm for basis extension (Lecture Week 4).}
To extend an independent set $S=\{v_1,\ldots,v_m\}$ to a basis of $V$ with $\dim(V)=n$:
\begin{enumerate}[leftmargin=*]
\item Append a known basis $\{b_1,\ldots,b_n\}$ of $V$ to form the list $v_1,\ldots,v_m,b_1,\ldots,b_n$.
\item Convert to coordinate vectors (columns in $\mathbb{F}^n$) if necessary.
\item Row reduce the combined matrix.  Identify which of the appended basis vectors $b_j$ correspond to non-pivot columns---these are the \emph{redundant} ones.
\item Remove the redundant $b_j$'s.  The remaining list $v_1,\ldots,v_m,b_{i_1},\ldots,b_{i_{n-m}}$ is a basis of $V$.
\end{enumerate}
As a bonus, $V=\operatorname{Span}(S)\oplus\operatorname{Span}\{b_{i_1},\ldots,b_{i_{n-m}}\}$, giving a direct-sum complement ``for free.''
\end{policynote}

\begin{theorem}{Spanning sets contain a basis (Tutorial 3, Question 3)}{ch2-spanning-contains-basis}

Let $V$ be finite-dimensional and let $S\subseteq V$ be a spanning set: $\operatorname{Span}(S)=V$.

Then $S$ contains a basis of $V$.

\begin{proof}

Start with a finite spanning subset of $S$ (in finite dimensions, spanning implies there is a finite spanning subset).

List it as $(u_1,\ldots,u_m)$.

If $\{u_1,\ldots,u_m\}$ is linearly independent, it is already a basis.

Otherwise it is dependent; by Theorem~\ref{thm:ch2-dep-iff-in-span}, one of the $u_k$ lies in the span of the others, so removing $u_k$ does not change the span.

Repeat until the remaining list is linearly independent. The process terminates because the list size strictly decreases.

The remaining set is linearly independent and still spans $V$, hence is a basis contained in $S$.

\end{proof}

\end{theorem}

% --- ADDED: Two complementary directions summary ---
\begin{keyformula}
\textbf{Two directions for building a basis.}
\begin{itemize}[leftmargin=*]
\item \emph{Bottom-up (extension):} start with an independent set, keep adding vectors outside the span.  Terminates when the set spans $V$.
\item \emph{Top-down (extraction):} start with a spanning set, keep removing vectors that lie in the span of the remaining ones.  Terminates when the set is independent.
\end{itemize}
Both produce a basis.  The choice depends on what you start with.
\end{keyformula}

\begin{examtip}

\begin{itemize}

\item \textbf{To prove a set is linearly independent:} start from $\sum \alpha_i v_i=0$ and force $\alpha_i=0$.

\item \textbf{To prove a set spans:} take an arbitrary $v$ and explicitly solve for coefficients.

\item \textbf{Fast finite-dimensional shortcut:} if you already know $\dim(V)=n$, then for $n$ vectors,

independence $\Longleftrightarrow$ spanning $\Longleftrightarrow$ basis.

\end{itemize}

\end{examtip}

%--------------------------------------------------------------------

\subsection{Dimension Computations and Dimension Formulae}

%--------------------------------------------------------------------

\begin{examplebox}

\textbf{Worked Example (Week 4 lecture): dimension of a plane in $\mathbb{R}^3$.}

Let

\[
\Pi=\left\{(x,y,z)\in\mathbb{R}^3 \ \middle|\ x+2y+3z=0\right\}.
\]

Find a basis and $\dim(\Pi)$.

\medskip

\textbf{Solution.}

Solve $x=-2y-3z$ and parametrize by $y=s$, $z=t$:

\[
(x,y,z)=(-2s-3t,\ s,\ t)=s(-2,1,0)+t(-3,0,1).
\]

Thus

\[
\Pi=\operatorname{Span}\{(-2,1,0),\ (-3,0,1)\}.
\]

The two spanning vectors are linearly independent, hence form a basis. Therefore $\dim(\Pi)=2$.

\end{examplebox}

% --- ADDED: General method for constraint-defined subspaces ---
\begin{policynote}
\textbf{General recipe for constraint-defined subspaces.}
If $W=\{v\in V\mid \text{linear conditions on }v\}$:
\begin{enumerate}[leftmargin=*]
\item Write out the constraints as linear equations in the coordinates/entries.
\item Solve the system: express the constrained variables in terms of the free variables.
\item Parametrize: write a generic element of $W$ as a linear combination of fixed vectors (one per free parameter).
\item The fixed vectors form a spanning set for $W$; check independence (usually immediate from the parametrisation).
\end{enumerate}
The number of free parameters equals $\dim(W)$.
As a sanity check: each \emph{independent} linear constraint reduces the dimension by $1$ (when the constraints are independent), so $\dim(W)\ge\dim(V)-(\text{number of constraints})$.
\end{policynote}

\begin{examplebox}

\textbf{Worked Example (Week 4 lecture): dimension of upper-triangular matrices.}

Let $U_n=\{M\in\mathbb{R}^{n,n}\mid M \text{ is upper triangular}\}$.

Find $\dim(U_n)$.

\medskip

\textbf{Solution.}

A matrix in $U_n$ has free entries in positions $(i,j)$ with $1\le i\le j\le n$.

There are

\[
1+2+\cdots+n=\frac{n(n+1)}{2}
\]

such positions. A natural basis is $\{E_{ij}\mid 1\le i\le j\le n\}$ (matrix units), so

\[
\dim(U_n)=\frac{n(n+1)}{2}.
\]

\end{examplebox}

\begin{theorem}{Dimension of a sum}{ch2-dim-sum}

Let $V$ be a finite-dimensional vector space and let $V_1,V_2$ be subspaces of $V$.

Then $V_1+V_2$ and $V_1\cap V_2$ are finite-dimensional and

\[
\dim(V_1+V_2)=\dim(V_1)+\dim(V_2)-\dim(V_1\cap V_2).
\]

\begin{proof}

Let $A$ be a basis of $V_1\cap V_2$.

Extend $A$ to a basis of $V_1$ by adding a set $B$ (so $A\cup B$ is a basis of $V_1$),

and extend $A$ to a basis of $V_2$ by adding a set $C$ (so $A\cup C$ is a basis of $V_2$).

One shows that $A\cup B\cup C$ is linearly independent and spans $V_1+V_2$, hence is a basis of $V_1+V_2$.

Therefore

\[
\dim(V_1+V_2)=|A|+|B|+|C|
=\bigl(|A|+|B|\bigr)+\bigl(|A|+|C|\bigr)-|A|
=\dim(V_1)+\dim(V_2)-\dim(V_1\cap V_2).
\]

\end{proof}

\end{theorem}

% --- ADDED: Full proof that A ∪ B ∪ C is independent ---
\begin{policynote}
\textbf{Why is $A\cup B\cup C$ linearly independent? (Lecture Week 4, full detail).}

Suppose $\sum_{a\in A}\lambda_a\,a + \sum_{b\in B}\mu_b\,b + \sum_{c\in C}\kappa_c\,c = 0$.
Then $\sum_{b\in B}\mu_b\,b = -\sum_{a\in A}\lambda_a\,a - \sum_{c\in C}\kappa_c\,c$.
The left side lies in $V_1$ (since $B\subseteq V_1$), and the right side can be rewritten as $\sum(-\lambda_a)\,a + \sum(-\kappa_c)\,c\in V_2$ (since $A\subseteq V_1\cap V_2\subseteq V_2$ and $C\subseteq V_2$).

Hence $\sum\mu_b\,b\in V_1\cap V_2=\operatorname{Span}(A)$.
But $\sum\mu_b\,b\in\operatorname{Span}(B)$, and $\operatorname{Span}(A)\cap\operatorname{Span}(B)=\{0\}$ (since $A\cup B$ is a basis of $V_1$, and by the Three-Concepts-in-One Theorem from Chapter~1).
So $\sum\mu_b\,b=0$, and by independence of $B$, all $\mu_b=0$.

Similarly, $\sum\kappa_c\,c\in V_1\cap V_2=\operatorname{Span}(A)$, and $\operatorname{Span}(A)\cap\operatorname{Span}(C)=\{0\}$, so all $\kappa_c=0$.

Substituting back: $\sum\lambda_a\,a=0$, so all $\lambda_a=0$ by independence of $A$.
Hence $A\cup B\cup C$ is linearly independent.
\end{policynote}

% --- ADDED: Analogy with inclusion-exclusion ---
\begin{policynote}
\textbf{Analogy with inclusion-exclusion (Lecture Week 4).}
The dimension formula $\dim(V_1+V_2)=\dim(V_1)+\dim(V_2)-\dim(V_1\cap V_2)$ is the vector-space analogue of the set-theoretic inclusion-exclusion formula $|A\cup B|=|A|+|B|-|A\cap B|$.

\emph{Caution (also from Lecture Week 4):} this analogy does \emph{not} extend to three subspaces.  The ``three-subspace inclusion-exclusion'' formula
$\dim(V_1+V_2+V_3)=\dim(V_1)+\dim(V_2)+\dim(V_3)-\dim(V_1\cap V_2)-\dim(V_2\cap V_3)-\dim(V_1\cap V_3)+\dim(V_1\cap V_2\cap V_3)$
is \emph{false} in general for vector subspaces.
\end{policynote}

\begin{keyformula}

\textbf{Dimension identity and immediate inequalities.}

For subspaces $U,V\subseteq \mathbb{F}^n$,

\[
\dim(U\cap V)=\dim(U)+\dim(V)-\dim(U+V).
\]

Since $\dim(U+V)\le n$, we get the useful bound

\[
\dim(U\cap V)\ \ge\ \dim(U)+\dim(V)-n.
\]

This is often the fastest way to prove intersections are nontrivial.

\end{keyformula}

\begin{examplebox}

\textbf{Worked Example (Tutorial 3, Question 7): two $5$D subspaces in $\mathbb{R}^9$ intersect nontrivially.}

Let $U$ and $V$ be $5$-dimensional subspaces of $\mathbb{R}^9$. Show $U\cap V\neq\{0\}$.

\medskip

\textbf{Solution.}

By Theorem~\ref{thm:ch2-dim-sum},

\[
\dim(U\cap V)=\dim(U)+\dim(V)-\dim(U+V)\ge 5+5-9=1,
\]

since $U+V\subseteq\mathbb{R}^9$ implies $\dim(U+V)\le 9$.

Thus $\dim(U\cap V)\ge 1$, so $U\cap V$ contains a nonzero vector.

\end{examplebox}

\begin{extension}[title={Extension (Tutorial 3, Question 7, Method 2)}]

If $U\cap V=\{0\}$, then $\mathbb{R}^9\supseteq U+V=U\oplus V$, so

\[
\dim(U+V)=\dim(U)+\dim(V)=10.
\]

But any subspace of $\mathbb{R}^9$ has dimension at most $9$, contradiction.

Hence $U\cap V\neq\{0\}$.

\end{extension}

% --- ADDED: Lecture Week 4 example on linear planes ---
\begin{examplebox}
\textbf{Worked Example (Week 4 lecture): intersection of two planes in $\mathbb{R}^3$.}

Let $\Pi_1$ and $\Pi_2$ be linear planes (i.e.\ $2$-dimensional subspaces) in $\mathbb{R}^3$.
What are the possible dimensions for $\Pi_1\cap\Pi_2$?

\medskip

\textbf{Solution.}
By the dimension formula,
\[
\dim(\Pi_1\cap\Pi_2)=\dim(\Pi_1)+\dim(\Pi_2)-\dim(\Pi_1+\Pi_2)=4-\dim(\Pi_1+\Pi_2).
\]
Since $\Pi_1+\Pi_2$ is a subspace of $\mathbb{R}^3$, we have $2\le\dim(\Pi_1+\Pi_2)\le 3$ (it must contain $\Pi_1$, so $\dim\ge 2$).

\begin{itemize}[leftmargin=*]
\item If $\dim(\Pi_1+\Pi_2)=3$: then $\dim(\Pi_1\cap\Pi_2)=1$ (a line through the origin).
\item If $\dim(\Pi_1+\Pi_2)=2$: then $\dim(\Pi_1\cap\Pi_2)=2$, which forces $\Pi_1=\Pi_2$ (same plane).
\end{itemize}
Both cases are realisable, so the possible dimensions are $\{1,2\}$.
\end{examplebox}

\begin{definition}{Direct sum (as used in Weeks 3--4)}{ch2-direct-sum}

Let $U,W$ be subspaces of $V$.

We write $V=U\oplus W$ if

\[
V=U+W
\quad\text{and}\quad
U\cap W=\{0\}.
\]

In this case, each $v\in V$ has a \emph{unique} decomposition $v=u+w$ with $u\in U$, $w\in W$.

\end{definition}

\begin{theorem}{Dimension of a direct sum}{ch2-dim-direct-sum}

If $V=U\oplus W$ with $V$ finite-dimensional, then

\[
\dim(V)=\dim(U)+\dim(W).
\]

\begin{proof}

Since $U\cap W=\{0\}$, Theorem~\ref{thm:ch2-dim-sum} gives

\[
\dim(U+W)=\dim(U)+\dim(W)-\dim(U\cap W)=\dim(U)+\dim(W).
\]

But $U+W=V$, so $\dim(V)=\dim(U)+\dim(W)$.

\end{proof}

\end{theorem}

% --- ADDED: Useful criterion for direct sum via dimensions ---
\begin{keyformula}
\textbf{Checking a direct sum via dimensions (very common on exams).}
Let $U,W\le V$ with $V$ finite-dimensional.  Then $V=U\oplus W$ if and only if \emph{any two} of the following three conditions hold:
\begin{enumerate}[leftmargin=*]
\item $V=U+W$;
\item $U\cap W=\{0\}$;
\item $\dim(U)+\dim(W)=\dim(V)$.
\end{enumerate}
\emph{Why any two suffice:} conditions (1) and (2) are the definition; if (1) and (3) hold, the dimension formula gives $\dim(U\cap W)=0$, so (2) follows; if (2) and (3) hold, the dimension formula gives $\dim(U+W)=\dim(V)$, so (1) follows.

This shortcut often saves a full spanning argument: just verify (2) and (3).
\end{keyformula}

\begin{examplebox}

\textbf{Worked Example (Week 4 lecture): a direct-sum conclusion from dimensions.}

Let $U,V\subseteq\mathbb{R}^8$ be subspaces such that $\dim(U)=3$, $\dim(V)=5$, and $\dim(U+V)=8$.

Show that $\mathbb{R}^8=U\oplus V$.

\medskip

\textbf{Solution.}

By Theorem~\ref{thm:ch2-dim-sum},

\[
\dim(U\cap V)=\dim(U)+\dim(V)-\dim(U+V)=3+5-8=0.
\]

Hence $U\cap V=\{0\}$. Since $U+V=\mathbb{R}^8$ by the dimension assumption, we conclude $\mathbb{R}^8=U\oplus V$.

\end{examplebox}

% --- ADDED: Week 4 lecture example on P_4(R) with kernel viewpoint ---
\begin{examplebox}
\textbf{Worked Example (Week 4 lecture): basis, extension, and direct-sum complement for $\{P\in P_4(\mathbb{R})\mid P(6)=0\}$.}

Define $V=\{P\in P_4(\mathbb{R})\mid P(6)=0\}$.

\medskip
\emph{(a) Basis and dimension.}
The evaluation map $\operatorname{ev}_6:P_4(\mathbb{R})\to\mathbb{R}$ given by $\operatorname{ev}_6(P)=P(6)$ is a nonzero linear functional, so $\dim(\ker(\operatorname{ev}_6))=5-1=4$.
Every polynomial vanishing at $6$ has $(t-6)$ as a factor (up to degree consideration): $P(t)=(t-6)Q(t)$ with $\deg(Q)\le 3$.
A basis is $\{t-6,\;t(t-6),\;t^2(t-6),\;t^3(t-6)\}$, confirming $\dim(V)=4$.

\medskip
\emph{(b) Extension and complement.}
Choose any $p\notin V$, e.g.\ $p=1$ since $1(6)=6\neq 0$.
Then $\{t-6,\;t(t-6),\;t^2(t-6),\;t^3(t-6),\;1\}$ is a basis of $P_4(\mathbb{R})$, and $W=\operatorname{Span}\{1\}$ satisfies $P_4(\mathbb{R})=V\oplus W$.
\end{examplebox}

%--------------------------------------------------------------------

\subsection{High-Yield Worked Problems from Tutorial 3}

%--------------------------------------------------------------------

\begin{examplebox}

\textbf{Worked Example (Tutorial 3, Question 4): commuting matrices form a subspace and its dimension.}

Let

\[
A=\begin{pmatrix}1&2\\0&1\end{pmatrix},
\qquad
\mathcal{V}=\left\{B\in\mathbb{R}^{2,2}\ \middle|\ AB=BA\right\}.
\]

Find $\dim(\mathcal{V})$.

\medskip

\textbf{Solution.}

Write $B=\begin{pmatrix}x&y\\ z&t\end{pmatrix}$.

Compute

\[
AB=\begin{pmatrix}x+2z&y+2t\\ z&t\end{pmatrix},
\qquad
BA=\begin{pmatrix}x&2x+y\\ z&2z+t\end{pmatrix}.
\]

Equating entries $AB=BA$ gives

\[
x+2z=x \Rightarrow z=0,
\qquad
y+2t=2x+y \Rightarrow t=x,
\]

and the remaining equations are automatic once $z=0$.

Thus

\[
B=\begin{pmatrix}x&y\\0&x\end{pmatrix}
=x\begin{pmatrix}1&0\\0&1\end{pmatrix}
+y\begin{pmatrix}0&1\\0&0\end{pmatrix}.
\]

Therefore $\mathcal{V}=\operatorname{Span}\!\left\{I,\ \begin{pmatrix}0&1\\0&0\end{pmatrix}\right\}$ and $\dim(\mathcal{V})=2$.

\end{examplebox}

\begin{extension}[title={Extension (Tutorial 3, Question 4, Method 2)}]

Write $A=I+2N$ with $N=\begin{pmatrix}0&1\\0&0\end{pmatrix}$ and $N^2=0$.

Since $I$ commutes with everything, $AB=BA$ iff $NB=BN$.

For $B=\begin{pmatrix}x&y\\ z&t\end{pmatrix}$,

\[
NB=\begin{pmatrix}z&t\\0&0\end{pmatrix},
\qquad
BN=\begin{pmatrix}0&x\\0&z\end{pmatrix}.
\]

Thus $NB=BN$ forces $z=0$ and $t=x$, giving $B=xI+yN$ and hence $\dim(\mathcal{V})=2$ with fewer computations.

\end{extension}

% --- ADDED: General commutant observation ---
\begin{policynote}
\textbf{Why the nilpotent decomposition helps (general principle).}
If $A=\alpha I+N$ where $N$ is nilpotent, then $AB=BA\iff NB=BN$ because $I$ commutes with everything.
This reduces the problem to commuting with a simpler (nilpotent) matrix.
For $2\times 2$ matrices, the commutant of a non-scalar matrix always has dimension $2$; the commutant of a scalar matrix ($A=\alpha I$) is all of $\mathbb{R}^{2,2}$, with dimension $4$.
\end{policynote}

\begin{examplebox}

\textbf{Worked Example (Tutorial 3, Question 5): a polynomial subspace, basis extension, and direct sum.}

Let

\[
V=\left\{P\in P_4(\mathbb{R})\ \middle|\ P(2)=P(5)\right\}.
\]

\begin{enumerate}

\item Find a basis of $V$ and $\dim(V)$.

\item Extend your basis to a basis of $P_4(\mathbb{R})$.

\item Find a subspace $W$ such that $P_4(\mathbb{R})=V\oplus W$.

\end{enumerate}

\medskip

\textbf{Solution.}

Let $q(t)=(t-2)(t-5)$ (degree $2$). For any $Q$, $(qQ)(2)=(qQ)(5)=0$, so $qQ\in V$.

Every $P\in P_4(\mathbb{R})$ can be written uniquely as

\[
P(t)=q(t)Q(t)+L(t),
\]

where $\deg(Q)\le 2$ and $\deg(L)\le 1$ (polynomial division by $q$). Write $L(t)=at+b$.

Then

\[
P(2)=L(2)=2a+b,\qquad P(5)=L(5)=5a+b,
\]

so the constraint $P(2)=P(5)$ forces $2a+b=5a+b$, hence $a=0$.

Thus $L$ must be constant, and we obtain the description

\[
V=\left\{q(t)Q(t)+c \ \middle|\ \deg(Q)\le 2,\ c\in\mathbb{R}\right\}.
\]

A convenient basis is

\[
\mathcal{B}_V=\left\{1,\ q(t),\ t\,q(t),\ t^2 q(t)\right\},
\]

so $\dim(V)=4$.

Since $\dim(P_4(\mathbb{R}))=5$, we extend by adding any polynomial not in $V$; for instance $t\notin V$ because $t(2)\neq t(5)$.

Hence

\[
\mathcal{B}=\left\{1,\ q(t),\ t\,q(t),\ t^2 q(t),\ t\right\}
\]

is a basis of $P_4(\mathbb{R})$.

Let $W=\operatorname{Span}\{t\}$. Then $V+W=P_4(\mathbb{R})$ by the basis above.

Also $V\cap W=\{0\}$: if $ct\in V$, then $(ct)(2)=(ct)(5)$ implies $2c=5c$, so $c=0$.

Therefore $P_4(\mathbb{R})=V\oplus W$.

\end{examplebox}

\begin{extension}[title={Extension (Tutorial 3, Question 5, Method 2)}]

Define the linear functional $L:P_4(\mathbb{R})\to\mathbb{R}$ by $L(P)=P(2)-P(5)$.

Then $V=\ker(L)$. Since $L\neq 0$, $\dim(\ker(L))=\dim(P_4)-1=4$.

To build a complement, choose any $p$ with $L(p)\neq 0$; for example $p(t)=t$ gives $L(t)=2-5=-3$.

Then $P_4(\mathbb{R})=\ker(L)\oplus \operatorname{Span}\{t\}$, with decomposition

$P=\bigl(P-\tfrac{L(P)}{L(t)}t\bigr)+\tfrac{L(P)}{L(t)}t$.

\end{extension}

% --- ADDED: General codimension-1 complement recipe ---
\begin{examtip}
\textbf{General recipe for complementing a codimension-$1$ subspace.}
If $V=\ker(L)$ for a nonzero linear functional $L:W\to\mathbb{F}$, then:
\begin{itemize}[leftmargin=*]
\item $\dim(V)=\dim(W)-1$.
\item Pick \emph{any} $p\in W$ with $L(p)\neq 0$.
\item Then $W=V\oplus\operatorname{Span}\{p\}$.
\item The explicit decomposition is: $w = \bigl(w-\tfrac{L(w)}{L(p)}p\bigr)+\tfrac{L(w)}{L(p)}p$.
\end{itemize}
This pattern appears repeatedly in MH1201 (trace-zero matrices, evaluation constraints, etc.).
\end{examtip}

\begin{examplebox}

\textbf{Worked Example (Week 4 lecture / dimension application): translate a large independent set.}

Let $v_1,\ldots,v_{2024}\in V$ be linearly independent. Show that for any $w\in V$,

\[
\dim\Bigl(\operatorname{Span}\{v_1+w,\ldots,v_{2024}+w\}\Bigr)\ge 2023.
\]

\medskip

\textbf{Solution.}

Consider the $2023$ vectors

\[
(v_1+w)-(v_2+w)=v_1-v_2,\ \ (v_2+w)-(v_3+w)=v_2-v_3,\ \ \ldots,\ \ (v_{2023}+w)-(v_{2024}+w)=v_{2023}-v_{2024}.
\]

Each difference lies in $\operatorname{Span}\{v_1+w,\ldots,v_{2024}+w\}$, so it suffices to show that

\[
\{v_1-v_2,\ v_2-v_3,\ \ldots,\ v_{2023}-v_{2024}\}
\]

is linearly independent.

Suppose $\sum_{i=1}^{2023}\alpha_i (v_i-v_{i+1})=0$. Expanding gives

\[
\alpha_1 v_1+(\alpha_2-\alpha_1)v_2+\cdots+(\alpha_{2023}-\alpha_{2022})v_{2023}-\alpha_{2023}v_{2024}=0.
\]

Since $\{v_1,\ldots,v_{2024}\}$ is linearly independent, all coefficients must be zero:

\[
\alpha_1=0,\ \ \alpha_2-\alpha_1=0,\ \ldots,\ \alpha_{2023}-\alpha_{2022}=0,\ \ -\alpha_{2023}=0.
\]

Hence $\alpha_1=\cdots=\alpha_{2023}=0$, proving the differences are independent.

Therefore the span of $\{v_1+w,\ldots,v_{2024}+w\}$ contains $2023$ independent vectors, so its dimension is at least $2023$.

\end{examplebox}

\begin{extension}[title={Extension (Tutorial 3, Question 6, Method 2)}]

\textbf{Stronger fact (and a structural proof).}

If $\{v_1,\ldots,v_n\}$ is linearly independent and $w\notin\operatorname{Span}\{v_1,\ldots,v_n\}$, then $\{v_1+w,\ldots,v_n+w\}$ is linearly independent.

Let $U=\operatorname{Span}\{v_1,\ldots,v_n,w\}$ and define $T:U\to U$ by $T(v_i)=v_i+w$ and $T(w)=w$, extended linearly.

Then $T$ is invertible with inverse $S(v_i)=v_i-w$, $S(w)=w$, so $T$ is injective; injective linear maps send independent sets to independent sets, hence $\{v_i+w\}$ is independent.

\end{extension}

% --- ADDED: Tutorial 3 Q8 on Q-vector spaces ---
\begin{examplebox}
\textbf{Worked Example (Tutorial 3, Question 8): $\mathbb{R}$ as a $\mathbb{Q}$-vector space.}

Let $\mathbb{R}_{\mathbb{Q}}$ denote $\mathbb{R}$ viewed as a vector space over $\mathbb{Q}$ (with usual addition and $\mathbb{Q}$-scalar multiplication).
For a fixed $r\in\mathbb{R}$, define $V=\{a+br\mid a,b\in\mathbb{Q}\}$.

\begin{enumerate}[leftmargin=*]
\item $V$ is a $\mathbb{Q}$-subspace of $\mathbb{R}_{\mathbb{Q}}$: it is non-empty ($0\in V$), closed under addition ($\mathbb{Q}$-sums of rationals are rational), and closed under $\mathbb{Q}$-scalar multiplication ($\mathbb{Q}\cdot\mathbb{Q}\subseteq\mathbb{Q}$).
\item If $r\in\mathbb{Q}$: then $V=\{a+br\mid a,b\in\mathbb{Q}\}=\mathbb{Q}$ (every element is rational), so $\dim_{\mathbb{Q}}(V)=1$ with basis $\{1\}$.
\item If $r\notin\mathbb{Q}$: then $\{1,r\}$ spans $V$ by definition.  For independence: if $\alpha\cdot 1+\beta\cdot r=0$ with $\alpha,\beta\in\mathbb{Q}$ and $\beta\neq 0$, then $r=-\alpha/\beta\in\mathbb{Q}$, contradiction.  Hence $\{1,r\}$ is a $\mathbb{Q}$-basis and $\dim_{\mathbb{Q}}(V)=2$.
\end{enumerate}

\emph{Key takeaway:} dimension depends on the field.  Over $\mathbb{R}$, $\dim_{\mathbb{R}}(\mathbb{R})=1$.  Over $\mathbb{Q}$, $\dim_{\mathbb{Q}}(\mathbb{R})$ is infinite (uncountable, in fact).
\end{examplebox}

\begin{commonmistake}

\begin{itemize}

\item Forgetting that \textbf{dimension depends on the field}. For example, $\mathbb{R}$ as a vector space over $\mathbb{Q}$ has very different dimension properties than over $\mathbb{R}$.

\item Misusing the dimension formula: it requires \emph{subspaces} of a \emph{finite-dimensional} vector space.

\end{itemize}

\end{commonmistake}

\begin{examtip}

\textbf{Dimension-proof checklist (very common in MH1201).}

\begin{itemize}

\item If asked to show $U\cap V\neq\{0\}$, try $\dim(U\cap V)\ge \dim(U)+\dim(V)-\dim(\text{ambient})$.

\item If given $\dim(U)+\dim(V)=\dim(U+V)$, then $\dim(U\cap V)=0$ and (if $U+V$ is the whole space) you get a direct sum.

\item If asked to build a basis for a constraint-defined subspace (e.g.\ $P(2)=P(5)$, $\operatorname{trace}(M)=0$, $AB=BA$):

translate the condition into \emph{linear equations in coefficients/entries}, then parametrize and read off a spanning set; finally check independence (often immediate).

\end{itemize}

\end{examtip}

%==================================================
% USER COMMENTS (EDITABLE)
%==================================================
% - (Write personal reminders, TODOs, or links here.)
% - TODO: Memorize the "any two of three" criterion for direct sums (sum, trivial intersection, dimension sum).
% - TODO: Practice the practical basis-extension algorithm (append + row reduce) on Tutorial 3 Q5 and Week 4 P(6)=0 example.
% - TODO: Know the standard dimension table by heart (F^n, P_n, F^{m,n}, trace-zero, symmetric, skew-symmetric, upper-triangular).
% - TODO: Re-derive the full proof that A ∪ B ∪ C is independent in the dimension-of-sum theorem.
% - TODO: Remember that the three-subspace inclusion-exclusion formula FAILS for vector spaces.
% - Midterm focus: Size-Limitation Lemma proof, Basis Extension, dimension formula applications.
